% Part: first-order-logic
% Chapter: natural-deduction
% Section: derivations

\documentclass[../../../include/open-logic-section]{subfiles}

\begin{document}

\iftag{FOL}
      {\olfileid[hi]{fol}{ntd}{der}}
      {\olfileid[hi]{pl}{ntd}{der}}

\olsection{\usetoken{p}{derivation}}

\begin{explain}
हम बता चुके हैं कि मान्यता क्या है और अनुमान-नियम भी दे चुके हैं। स्वाभाविक
निगमन की !!^{derivation}s इन्हीं से आगमनात्मक रूप से बनती हैं: हर
!!{derivation} या तो अकेली मान्यता है, या उसमें एक, दो अथवा तीन
!!{derivation}s के बाद एक सही अनुमान होता है।
\end{explain}

\begin{defn}[!!^{derivation}]
\Article{derivation} \emph{!!{derivation}}, !!a{sentence}~$!A$ को
मान्यताओं~$\Gamma$ से देने वाला !!{sentence}s का ऐसा परिमित वृक्ष है जो
निम्नलिखित शर्तें पूरी करे:
\begin{enumerate}
\item वृक्ष के शीर्षस्थ !!{sentence}s या तो $\Gamma$ में हैं या वृक्ष के किसी
  अनुमान द्वारा !!{discharged} किए जाते हैं।
\item वृक्ष का सबसे निचला !!{sentence}~$!A$ है।
\item वृक्ष का हर !!{sentence}, सबसे नीचे के~$!A$ को छोड़कर, किसी अनुमान-नियम
  के सही प्रयोग का पूर्वाधार है; उस प्रयोग का निष्कर्ष वृक्ष में उस
  !!{sentence} के ठीक नीचे स्थित है।
\end{enumerate}
तब $!A$ को !!{derivation} का \emph{निष्कर्ष} और $\Gamma$ को उसकी
!!{undischarged} मान्यताएँ कहते हैं।

यदि !!a{derivation} ऐसी हो जो $!A$ को~$\Gamma$ से देती हो, तो कहते हैं कि
$!A$,~$\Gamma$ से \emph{!!{derivable}} है, या प्रतीकों में: $\Gamma \Proves
!A$। यदि ऐसी !!a{derivation} हो जिसमें~$!A$ निष्कर्ष हो और हर मान्यता
!!{discharged} हो, तो~$\Proves !A$ लिखते हैं।
\end{defn}

\begin{ex}
हर मान्यता स्वयं !!a{derivation} है; अतः $!A$ स्वयं !!a{derivation} है और
$!B$ भी। इन पर $\Intro{\land}$ नियम लगाने से नई !!{derivation} मिलती है:
\begin{prooftree}
\AxiomC{$!A$}
\AxiomC{$!B$}
\RightLabel{\Intro{\land}}
\BinaryInfC{$!A \land !B$}
\end{prooftree}
ये नियम सामान्य हैं: इनमें $!A$ और~$!B$ को किन्हीं भी !!{sentence}s, जैसे $!C$
और~$!D$, से बदल सकते हैं। तब निष्कर्ष $!C \land !D$ होगा; अतः
\begin{prooftree}
\AxiomC{$!C$}
\AxiomC{$!D$}
\RightLabel{\Intro{\land}}
\BinaryInfC{$!C \land !D$}
\end{prooftree}
सही !!{derivation} है। मान्यताएँ उलटकर $!D$ को~$!A$ और $!C$ को~$!B$ की
भूमिका भी दे सकते हैं। अतः
\begin{prooftree}
\AxiomC{$!D$}
\AxiomC{$!C$}
\RightLabel{\Intro{\land}}
\BinaryInfC{$!D \land !C$}
\end{prooftree}
भी सही !!{derivation} है।

अब $\Intro{\lif}$ जैसा दूसरा नियम लगा सकते हैं। यह आपादन निष्कर्षित करने और
उसके पूर्वपद के समान किसी भी मान्यता को !!{discharge} करने देता है। अतः निम्न
दोनों सही !!{derivation}s हैं:
\begin{prooftree}
\AxiomC{$\Discharge{!C}{1}$}
\AxiomC{$!D$}
\RightLabel{\Intro{\land}}
\BinaryInfC{$!C \land !D$}
\DischargeRule{\Intro{\lif}}{1}
\UnaryInfC{$!C \lif (!C \land !D)$}
\DisplayProof\bottomAlignProof
\AxiomC{$!C$}
\AxiomC{$\Discharge{!D}{1}$}
\RightLabel{\Intro{\land}}
\BinaryInfC{$!C \land !D$}
\DischargeRule{\Intro{\lif}}{1}
\UnaryInfC{$!D \lif (!C \land !D)$}
\end{prooftree}
ये क्रमशः दिखाती हैं कि $!D \Proves !C \lif (!C \land !D)$ और
$!C \Proves !D \lif (!C \land !D)$।

याद रखें: मान्यताओं को मुक्त करना अनुमति है, अनिवार्यता नहीं। विशेषतः,
मान्यताएँ !!{derivation} में मौजूद न हों, तब भी नियम लगा सकते हैं। निम्नलिखित
अनुमत है, यद्यपि कोई मान्यता~$!A$ नहीं है जिसे !!{discharged} किया जाए:
\begin{prooftree}
  \AxiomC{$!B$}
  \DischargeRule{\Intro{\lif}}{1}
  \UnaryInfC{$!A \lif !B$}
\end{prooftree}
\end{ex}

\end{document}
